
Weight-space learning --- predicting neural network properties directly from their parameters --- faces a fundamental challenge: permuting a network's neurons produces a functionally identical model with different raw weights. Non-invariant encoders treat these equivalent configurations as distinct, and the cost is measurable: a sinusoidal positional encoder (CISE), designed as a representative non-invariant baseline, achieves $R^2 = -1.63$ when its predictions are averaged over functionally equivalent weight permutations, revealing that channel-position information is a primary predictive signal rather than background noise. This paper studies this failure through a three-step causal framework --- architecture determines within-orbit representational variance (OrbitVar), OrbitVar propagates to permutation-induced prediction error (MSE\_perm), and eliminating MSE\_perm improves downstream $R^2$ --- and measures each step empirically on ModelZooDataset CIFAR10-GS. Architecturally invariant encoders (DeepSets doubly-invariant sum pooling) achieve OrbitVar at machine precision (1.002e-14, twelve orders of magnitude below CISE's 0.010333), confirmed across all 100 models (Wilcoxon signed-rank p = 1.95e-18). For CISE, MSE\_perm/MSE\_total = 3.3452, demonstrating that permutation-induced prediction variance exceeds total prediction error by a factor of 3.35. DeepSets achieves $R^2$ = 0.9148 compared to CISE's $R^2$ = 0.8511, a 6.37 percentage-point improvement on the same testset with no additional training supervision. A pre-registered additive closure prediction ($\Delta MSE \approx MSE_perm^{C1}$) fails empirically (closure = 0.872), revealing an entanglement between MSE\_perm and the residual error component that opens questions about the geometry of non-invariant weight-space representations.

